%!TEX root = ./../../springer-mono-format/book.tex
%!TEX root = ./../../thesis-trier.tex

Here, we discuss the classes, and related literature, 
introduced in this chapter. 
Later, in \cref{sync_max:sec:subclasses,constr_sync:chap:constr_input,chap:prelim_complexity_results},
we investigate these subclasses w.r.t.~to decision problems, automata for sets of synchronizing
words, and the constrained synchronization problem, respectively.
A discussion of the results from the these chapters 
can be found in \cref{sync_max:sec:subclasses:discussion,constr_input:sec:discussion,prelim:compl:sec:discussion}.
%But we do not discuss the work of the author presented in this thesis related to these classes, as basically the whole thesis discusses this work.
Here, we leave out a discussion of synchronizable automata, because they are the topic of the whole work
and some related literature was already discussed in the introductory part (\cref{intro:chap:short_intro_sync_aut}).
The new notion of  $\rho$-automata is introduced here for the first time.
%Also, we do not discuss $\rho$-automata, as this type of automata is a new notion introduced here for the first time.

\paragraph{SI-Automata.} Rystsov introduced SI-automata in~\cite{Rystsov2000}
and proved the bound $2(n-1)^2$ for the reset length 
of an $n$-state SI-automata~$\mathcal A$.
In~\cite{Rystsov2022} the author showed the bound $n - 1$ in case the group of units of~$\TS{\mathcal A}$ is trivial and the bound
$(n-1)^2$ in case the group of units of~$\TS{\mathcal A}$ is a regular 
permutation group~\cite{Rystsov2022}. We deduce both bounds 
in \cref{sync_max:cor:augmentative_subset_IG,syncmax:graph:cor:OI-sync} in the case of OI-automata.

\paragraph{Permutation-Reset Automata.} Permutation-reset
automata were firstly mentioned by Zeiger
in the conference contribution~\cite{DBLP:conf/focs/Zeiger65} (see~\cite{DBLP:journals/iandc/Zeiger67a} for the extended version) 
in connection with the Krohn-Rhodes decomposition theory~\cite{KrohnRhodes65} (see \SEC~\ref{prelim:prelim:discussion}
for some background on this theory).
McNaughton~\cite{McNaughton67}
and Thierrin~\cite{DBLP:journals/mst/Thierrin68}
investigated permutation automata as language acceptors,
see also~\cite{Eilenberg1976,Pin86,DBLP:reference/hfl/Pin97} for information on the class of languages
accepted by permutation automata, the
so-called \emph{variety of group languages}. 
\IN{variety of group languages}

\paragraph{Commutative Automata.} Commutative automata
were introduced, to the best of our knowledge,
by Laing and Wright~\cite{laingwright62}.
They focused on commutative automata over a binary alphabet, gave
normal form theorems for regular expressions modulo commutativity,
and also considered infinite-state commutative automata
and regularity conditions. 
Rystsov~\cite{Rys96} gave (by an algebraic proof) the tight bound $n-1$
for a shortest synchronizing word of an $n$-state commutative automaton.
Fernau and Hoffmann~\cite{FernauHoffmann19} gave a different (combinatorial)
proof and a structure theorem for commutative automata that
reach this bound.
Fleck~\cite{DBLP:journals/jacm/Fleck62} determined the automorphism group
of strongly connected and commutative automata.
Cascade product decomposition results (see the discussion in~\cref{prelim:prelim:discussion} on automata theory for this notion) for commutative automata
were given in~\cite{Gelderie2010,DBLP:conf/lata/Gelderie11}.
See also~\cite{DBLP:journals/actaC/EsikI81a,DBLP:journals/actaC/EsikI81,DBLP:journals/actaC/Imreh81,Peak65} for further algebraic investigations.
Question on automata sizes for commutative regular languages
were investigated in~\cite{DBLP:journals/iandc/Hoffmann24,DBLP:journals/ijfcs/Hoffmann23a}.
The languages accepted by different types of commutative automata
were investigated in~\cite{DBLP:journals/actaC/AlmeidaEP17,
	DBLP:conf/cai/Hoffmann19,
	DBLP:conf/wia/Hoffmann21,
	Hoffmann23a,
	DBLP:journals/iandc/Hoffmann24,
	Pin86,
	DBLP:reference/hfl/Pin97}.
Algorithmic problems on commutative automata were investigated in~\cite{DBLP:journals/corr/abs-2108-05244},
and the membership problem on commutative transformation semigroups in~\cite{DBLP:journals/iandc/Beaudry88}.
 
\paragraph{Partially Ordered Automata.}
This type of automata has been considered in the literature
under different names: \emph{weakly acyclic}~\cite{Ryzhikov19a}, \emph{acyclic}~\cite{DBLP:conf/wia/JiraskovaM12},
\emph{extensive}~\cite{Pin2022}
or \emph{1-weak/very weak/linear}~\cite{DBLP:conf/lpar/DaxK08} (linear stems from the existence of a topological sorting of the states).
Ryzhikov~\cite{Ryzhikov19a} investigated several synchronization
and rank problems on partially ordered automata
and obtained the tight bound $n - r$ for a shortest word of rank at most~$r$.
Bruchertseifer and Fernau~\cite{DBLP:journals/ita/BruchertseiferF21}
investigated the subclass of \emph{series-parallel automata with loops}
with respect to questions of synchronizability in the setting of
parameterized complexity (see the part on parameterized complexity in~\cref{prelim:compl:sec:discussion} for more information). 
\mbox{Brzozowski} and Fich~\cite{DBLP:journals/jcss/BrzozowskiF80}
obtained that partially ordered automata
accept precisely those languages with $\mathcal R$-trivial syntactic monoids.
Jir\'askov\'a and Masopust~\cite{DBLP:conf/wia/JiraskovaM12} 
investigated the state complexity of the reverse operation for partially ordered automata,
Dax and Klaedtke~\cite{DBLP:conf/lpar/DaxK08} investigated
the translation of alternating partially ordered automata to nondeterministic automata
in the setting of infinite words.

\label{discussion:monotonic}
\paragraph{Monotonic Automata.} Ananichev and Volkov~\cite{DBLP:journals/tcs/AnanichevV04} introduced monotonic automata
and proved \cref{prelim:prop:monotonic_rank_r}.\footnote{The monotonic automata as introduced earlier 
	by Eppstein~\cite{Eppstein90} are more general, as they are only required to respect a cyclic order and were
	renamed \emph{oriented} automata by Ananichev and Volkov~\cite{DBLP:journals/tcs/AnanichevV04},
and \emph{cyclically monotonic} by Martyugin~\cite{Mar2009b}). Linearly and circularly monotonic might be another
	possible renaming for both concepts.}
However, transformation semigroups whose elements preserve some linear order
were investigated before, for example, in~\cite{Howie1971}.
The bound $n - 1$ for shortest synchronizing of monotonic automata
was stated in~\cite{AnanichevVolkov04} in the language of transformation semigroups.
In~\cite{RyzhikovS18} problems related to (subset) synchronization problems
on monotonic partial automata were investigated.
% (note that here we only consider
%complete automata as inputs to \constrsync{L}). \todo{complete oben definiert? sonst sagen transition fully defined.}
%komischer satz, konkret sagen welches problem
%An open problem
%from~\cite{DBLP:journals/tcs/AnanichevV04} was investigated in~\cite{Shcherbak05}.
Checking whether a given complete automaton is monotonic is \NP-complete~\cite{Szykula15}.
Anachichev and Volkov~\cite{AnanichevV05}
extended the tight bound $n-r$ on shortest words of rank at most~$r$ (\cref{prelim:prop:monotonic_rank_r})
to so-called \emph{generalized monotonic automata}
(these automata where defined in levels 
that yield an infinite hierarchy even in the case of strongly strongly connected 
automata~\cite{AlmeidaCardoso13})
Apart from its rather simple definition,
the behavior of monotonic automata as language acceptors 
remains rather elusive. Shcherbak~\cite{Shcherbak05} posed
the \emph{open problem} to give a combinatorial description (for example, by
regular expressions) of the languages accepted by monotonic automata.
In~\cite{DBLP:conf/icalp/Kopczynski06,Kopczynski08} monotonic
automata were used to describe certain winning conditions in infinite games.
In~\cite{DBLP:conf/dlt/Ananichev05} shortest words that are undefined on each state
in partially ordered monotonic automata were investigated,
which correspond to shortest synchronizing words in the complete automaton
resulting from adding a trap state.

\paragraph{Solvable Automata.} Rystsov~\cite{Rys96} gave the tight bound $n-1$
for shortest synchronizing words of $n$-state $R_1$- and $R_2$-solvable automata. Here, in \cref{prelim:thm:solvrank_r}, 
for solvable automata, we gave the bound $n - r$
for shortest words of rank at most $r$.
 In personal communication, Mikhail V. Volkov
 noted  that the transition monoid of a solvable automaton $\semiautomat{A}{Q}{\Sigma}{\delta}$ 
 is in the pseudovariety $\mathbf{DS}$ (in fact, of any automaton whose strongly connected components are blocks).
 Automata with this property
 admit a word of rank~$r$ and length at most~$n - r$
 if the minimal rank among its words is~$r$~\cite[Thm.~2.6]{AlmeidaS09}.
 \Cref{prelim:thm:solvrank_r}
 is slightly more general since it applies
 to values of~$r$ that do not equal
 the minimal rank among all the ranks of the words.
 The containment of their transition monoids in $\mathbf{DS}$
 follows from a characterization of $\mathbf{DS}$
 as the set of finite semigroups that are semilattices over semigroups
 nilpotent over their kernel (see the equivalence (4c) $\Leftrightarrow$ (1b) in~\cite[Thm.~3]{She95a}
 and~\cite{DBLP:journals/corr/abs-2409-06971,She95a} 
 for definitions of the mentioned notions).
 For let $P_1, \ldots, P_m$
 be a topological sorting of the strongly connected components w.r.t.~$\lereach$.
 Then set
 \[
 Y = \{ (x_1, \ldots, x_m) \in \{0,1\}^m \mid \exists w \in \Sigma^* \colon x_i = 1 \Leftrightarrow \delta(P_i,w) \subseteq P_i \}
 \]
 and, for $(x_1, \ldots, x_m) \in Y$, set
 \[
 S_{(x_1, \ldots, x_m)} = \{ \fword{\delta}{w} \colon Q \to Q \mid \forall i \in \{1,\ldots,m\} \colon x_i = 1 \Leftrightarrow \delta(P_i, w) \subseteq P_i \}.
 \]
 Then $Y$ is a semilattice with the operation
 \[
 (x_1, \ldots, x_m)(y_1, \ldots, y_m) = (\min\{x_1,y_1\},\ldots,\min\{x_m,y_m\})
 \]
 and $\{ S_y \}_{y \in Y}$
 forms a semilattice of semigroups
 where the $S_y$ are groups, and so nilpotent over their kernel.